% GENERATED FILE. Edit canonical lessons, then run npm run generate:books.
% canonical-source-hash: fnv1a64:15d92ecceca221b8
% canonical-lessons: LA-C46-lux, LA-C46-umbra, LA-C46-stella, LA-C46-luna, LA-C46-somnus

\chapter{Light, Shadow, and Going to Sleep}
\label{ch:la-lux}
\hlchaptermodality{46}

\begin{chapteropening}
\textbf{By the end of this chapter:} \emph{I can name light, shadow, star, moon and sleep in Latin, and use them to close a day the way a Latin speaker would.}

The fifth word of the group, taken apart against the four before it.
\end{chapteropening}

\section[lūx, lūcis]{lūx — the root of lucid, translucent and Lucifer, and a cousin of English light itself}

\label{lesson:LA-C46-lux}



\begin{quote}
A new group of five words begins here. Nothing is assumed but the alphabet you already read.
\end{quote}

\begin{cousinweb}
\textbf{lūx} — \textquotedblleft{}\textbf{light}.\textquotedblright{}

Proto-Indo-European \emph{\textbf{*lewk-}} is one of the best-attested roots there is, and it reaches English from three directions at once.

\textbf{Inherited:} English \textbf{light} is the Germanic descendant. Same ancestor, no borrowing.

\textbf{Borrowed from Latin:} something \textbf{lucid} is clear as light; something \textbf{translucent} lets light through; to \textbf{elucidate} is to shed light on. And \textbf{Lucifer} is \emph{lūx} plus \emph{ferre} — the \textbf{light-bringer}, which was the Romans' name for the morning star long before it was anything else.

\textbf{Borrowed from Greek:} the same root gave Greek \emph{leukos}, \textquotedblleft{}white,\textquotedblright{} which English carries in \textbf{leukemia}, named for white blood cells.

Three roads, one root — and the everyday word is the inherited one, which is the usual pattern in English rather than an accident here.
\end{cousinweb}

\subsection*{Your turn}
Say these aloud:

\begin{itemize}
\raggedright
  \item \emph{lūx, lūcis}
  \item every word this chapter has given you so far, in order
\end{itemize}

\begin{tcolorbox}[breakable,colback=teal!4,colframe=teal!35!black,title={Before you move on}]
What does \textbf{Lucifer} literally mean? (\textbf{Light-bringer} — the morning star.) Is English \emph{light} inherited or borrowed? (\textbf{Inherited}.) And which Greek word for \emph{white} is the same root? (\emph{\textbf{Leukos}} — leukemia.)
\end{tcolorbox}
\section[umbra, umbrae]{umbra — and the reason a little shade is called an umbrella}

\label{lesson:LA-C46-umbra}



\begin{quote}
Recall the previous word in this chapter before adding another.
\end{quote}

\begin{cousinweb}
\textbf{umbra} — \textquotedblleft{}\textbf{shade},\textquotedblright{} or a \textbf{shadow}.

The English word that gives it away is \textbf{umbrella} — Italian \emph{ombrella}, a \textbf{little shade}. It was a sunshade first and a rain shield only later, which is why a language of northern rain borrowed a word about southern sun.

To take \textbf{umbrage} is to take shade — offence, in the sense of being overshadowed. And \textbf{sombre} is \emph{sub} plus \emph{umbra} worn down through French: \textbf{under shadow}.

Astronomers kept the Latin bare. In an eclipse the \textbf{umbra} is the full shadow and the \textbf{penumbra} — \emph{paene}, \textquotedblleft{}almost\textquotedblright{} — is the partial one at its edge. That prefix is worth having: a \textbf{peninsula} is \emph{paene} plus \emph{īnsula}, an \textbf{almost- island}.

You have already met the \emph{shade} of evening under a different word, when this book taught you \emph{vesper}.
\end{cousinweb}

\subsection*{Your turn}
Say these aloud:

\begin{itemize}
\raggedright
  \item \emph{umbra, umbrae}
  \item every word this chapter has given you so far, in order
\end{itemize}

\begin{tcolorbox}[breakable,colback=teal!4,colframe=teal!35!black,title={Before you move on}]
What does \textbf{umbrella} literally mean? (\textbf{A little shade}.) What was it for first? (\textbf{Sun}, not rain.) And what does \emph{paene} add in \textbf{penumbra} and \textbf{peninsula}? (\textbf{Almost}.)
\end{tcolorbox}
\section[stēlla, stēllae]{stēlla — a cousin of English star, and the word inside constellation and disaster}

\label{lesson:LA-C46-stella}



\begin{quote}
Recall the previous word in this chapter before adding another.
\end{quote}

\begin{cousinweb}
\textbf{stēlla} — \textquotedblleft{}\textbf{a star}.\textquotedblright{}

A cousin again. Proto-Indo-European \emph{\textbf{*h\textsubscript{2}st\'{\={e}}r}} gave Latin \emph{stēlla}, Greek \emph{astēr}, and — by inheritance down the Germanic line — English \textbf{star}. Three branches, one ancestor.

From the Latin: something \textbf{stellar} is of the stars, and a \textbf{constellation} is a set of them standing \textbf{together} (\emph{con-} plus \emph{stēlla}).

From the Greek: an \textbf{asterisk} is a \textbf{little star}, \textbf{astronomy} is the naming of stars, and an \textbf{asteroid} is a star-shaped thing.

And one word that carries a whole worldview. \textbf{Disaster} is Italian \emph{disastro} — \emph{dis-}, \textquotedblleft{}ill,\textquotedblright{} plus \emph{astro}, \textquotedblleft{}star.\textquotedblright{} A \textbf{bad star}. The word only makes sense in a world where the stars decide what happens to you, and English has kept using it for four hundred years after it stopped believing that.
\end{cousinweb}

\subsection*{Your turn}
Say these aloud:

\begin{itemize}
\raggedright
  \item \emph{stēlla, stēllae}
  \item every word this chapter has given you so far, in order
\end{itemize}

\begin{tcolorbox}[breakable,colback=teal!4,colframe=teal!35!black,title={Before you move on}]
Is English \emph{star} a cousin or a loan? (\textbf{A cousin}.) What does \textbf{constellation} literally say? (\textbf{Stars standing together}.) And what does \textbf{disaster} literally say? (An \textbf{ill star}.)
\end{tcolorbox}
\section[lūna, lūnae]{lūna — the same root as lux, worn down to mean 'the shining one'}

\label{lesson:LA-C46-luna}



\begin{quote}
Recall the previous word in this chapter before adding another.
\end{quote}

\begin{cousinweb}
\textbf{lūna} — \textquotedblleft{}\textbf{the moon}.\textquotedblright{}

Look at what it is made of, because you have just met the pieces. \emph{Lūna} is \emph{\textbf{*lewk-snā}}: the \textbf{shining one}. It is the same root as \emph{lūx}, worn down until the join stopped showing.

You already know one place it lives. The Latin name for Monday is \emph{\textbf{diēs Lūnae}}, the moon's day — which this book taught you before you could take it apart, and which is exactly what German \emph{Montag} and English \emph{Monday} say too, in their own words.

\textbf{Lunar} is the plain loan. \textbf{Lunatic} is the interesting one: Latin \emph{lūnāticus} meant \textquotedblleft{}moon-struck,\textquotedblright{} from a belief that the moon's phases moved the mind. English has kept the word and dropped the theory, the same way it kept \emph{disaster}.

English \textbf{moon} is the inherited Germanic cousin — not from \emph{lūna}, but from a root meaning \textbf{measure}, because the moon is what a month is counted by.
\end{cousinweb}

\subsection*{Your turn}
Say these aloud:

\begin{itemize}
\raggedright
  \item \emph{lūna, lūnae}
  \item every word this chapter has given you so far, in order
\end{itemize}

\begin{tcolorbox}[breakable,colback=teal!4,colframe=teal!35!black,title={Before you move on}]
What is \emph{lūna} built from? (\emph{\textbf{Lūx}} — the shining one.) Which weekday is \emph{diēs Lūnae}? (\textbf{Monday}.) And what did \textbf{lunatic} originally claim? (\textbf{Moon-struck} — the moon moved the mind.)
\end{tcolorbox}
\section[somnus, somnī]{somnus — the word under insomnia and somnolent, and a cousin of Sanskrit svapna}

\label{lesson:LA-C46-somnus}



\begin{quote}
Recall the previous word in this chapter before adding another.
\end{quote}

\begin{cousinweb}
\textbf{somnus} — \textquotedblleft{}\textbf{sleep}.\textquotedblright{}

The last word of this chapter, and the right one to end a day on.

The root is Proto-Indo-European \emph{\textbf{*swep-}}, \textquotedblleft{}to sleep,\textquotedblright{} and it reaches a long way: Sanskrit \textbf{svapna}, \textquotedblleft{}sleep\textquotedblright{} or \textquotedblleft{}dream,\textquotedblright{} is the same word, and Old English had \emph{sweven} for a dream — a word that died out but is still in Chaucer.

\textbf{Insomnia} is \emph{in-}, \textquotedblleft{}not,\textquotedblright{} plus \emph{somnus}: \textbf{no sleep}. Someone \textbf{somnolent} is sleepy, and a \textbf{somnambulist} walks (\emph{ambulāre} — you have that verb) in their sleep.

The Roman god of sleep is \textbf{Somnus}, and his Greek counterpart \textbf{Hypnos} gave English \textbf{hypnosis}. Two languages, two gods, one job.

Now put the chapter together. \textbf{Lūx} goes, \textbf{umbra} comes, the \textbf{stēllae} and the \textbf{lūna} appear, and then \textbf{somnus}. That is a Latin evening in five words.
\end{cousinweb}

\subsection*{Your turn}
Say these aloud:

\begin{itemize}
\raggedright
  \item \emph{somnus, somnī}
  \item every word this chapter has given you so far, in order
\end{itemize}

\begin{tcolorbox}[breakable,colback=teal!4,colframe=teal!35!black,title={Before you move on}]
What does \textbf{insomnia} literally say? (\textbf{No sleep}.) What does a \textbf{somnambulist} do? (\textbf{Walks in their sleep} — \emph{ambulāre}.) And which Sanskrit word is \emph{somnus}'s cousin? (\emph{\textbf{Svapna}}.)
\end{tcolorbox}
